\section{Development of an integral-remainder quantum pricing method}
\label{sec:integral-remainder-quantum}

My research objective is to develop a distinct or improved quantum pricing
method for derivatives under an established classical financial model.
The underlying multivariate GBM remains unchanged. The preceding classical
certificates are enabling components and competitors; they are not a
substitute for the quantum construction. I now specify a candidate oracle
whose financial correction is represented by a bounded integral rather
than by subtracting nearly equal prices. The complete financial arithmetic
and its competitive gate budget remain an implementation task. I distinguish
the proved interface and realised auxiliary circuit from that unfinished
end-to-end comparison.

\subsection{The classical no-arbitrage target}

Let the physical stock system be
$dS_i/S_i=b_i\,dt+\Sigma_i\,dW$, with constant dividend yields $d_i$,
constant risk-free rate $r_f$ and a filtration generated by a
$d_B$-dimensional Brownian motion. Suppose $\Sigma\in\mathbb R^{n\times d_B}$
has full row rank. Choose a constant vector $\theta$ solving
$\Sigma\theta=b-r_f\mathbf1+d$; full row rank makes this equation solvable.
For horizon $T_{\rm mat}$ define
$Z_t=\exp(-\theta\cdot W_t-\|\theta\|^2t/2)$.
Independent normal increments and their moment generating function give
$\mathbb E[Z_t\mid\mathcal F_s]=Z_s$. In particular $Z_t>0$ and
$\mathbb E Z_{T_{\rm mat}}=1$, so $Z_{T_{\rm mat}}$ defines an equivalent
probability $Q$. Multiplication of every normal increment density by its
exponential factor shifts its mean by $-\theta\Delta t$ and preserves its
covariance and independence. Thus $W^Q_t=W_t+\theta t$ is Brownian under
$Q$ and the ex-dividend stock drift is $r_f-d_i$.

The discounted cumulative-gain process is
\[
 e^{-r_f t}S_i(t)+\int_0^t e^{-r_f u}d_iS_i(u)\,du.
\]
It\^o's formula gives zero drift and diffusion
$e^{-r_f t}S_i(t)\Sigma_i\,dW^Q_t$.
The explicit lognormal solution has finite second moments on this finite
horizon; hence the stochastic integral is square-integrable and the gain
is a true martingale. A self-financing admissible portfolio with zero
initial discounted wealth and wealth bounded below has a discounted
local-martingale gain. Stop a localising sequence, add its fixed lower-bound
constant and apply conditional Fatou's lemma. Its expectation at maturity
is at most zero. Nonnegative terminal wealth that is strictly positive
with positive probability contradicts this inequality, establishing
absence of such arbitrage.

I retain the notation of equation~\eqref{eq:target}:
$V=\mathbb E_Q F_H(A_{T_{\rm mat}})$ is the undiscounted mean and
$\Pi=e^{-r_fT_{\rm mat}}V$ is the time-zero derivative price.
At intermediate times the proposed derivative price is
$e^{-r_f(T_{\rm mat}-t)}\mathbb E_Q[F_H(A_{T_{\rm mat}})\mid\mathcal F_t]$.
Its discounted value is a bounded martingale. Augmenting the traded market
with this price process therefore preserves the preceding no-arbitrage
argument for admissible strategies. The construction is a model-consistent
price, rather than an assertion that a historical covariance calibration
identifies market option prices.

The independent finite-interface audit checks full volatility row rank and
the exact risk-neutral log drift in all nine frozen parameter models.
The price functional satisfies the no-arbitrage conditions just derived;
a finite estimate has the stated certified interval once its complete
evaluator, precision and confidence conditions have been met.
Separate rounded point estimates across strikes need not preserve every
joint static-arbitrage inequality. A claim about an entire tradable quoted
surface would require joint feasible prices or intervals, which is outside
this single-price accuracy result.

Full row rank is a sufficient condition here, not a necessity claim.
If $n=d_B$ and $\Sigma$ is invertible, its Brownian market-price-of-risk
vector is unique. If $d_B>n$, different vectors in $\ker\Sigma$ can give
different equivalent measures. In the constant-coefficient GBM considered
here they induce the same stock drift and covariance, so a terminal stock
basket has the same risk-neutral distribution under these choices.
Unspanned Brownian risks alone do not create price ambiguity for this
particular payoff. More general incomplete models require a stated choice
of $Q$. Quantum preparation changes the numerical estimator, not the
financial law or this no-arbitrage price functional.

\subsection{A bounded oracle for the exact Taylor correction}

Write $U(y)=T(y)-1\geq0$, $m=m_G(y)$ and let $r\geq1$ be fixed.
The exact order-$r-1$ pointwise correction is
\[
 R_r(y)=g_m(1+U)-\sum_{j=0}^{r-1}g_m^{(j)}(1)U^j/j!.
\]
Taylor's integral formula and the substitution $t=1+zU$ give, including
the continuous extension at $U=0$,
\begin{equation}
 R_r(y)={U^r\over r!}\mathbb E_Z[g_m^{(r)}(1+ZU)],\qquad
       Z\sim\operatorname{Beta}(1,r).
 \label{eq:quantum-beta-remainder}
\end{equation}
Indeed the integral remainder is
$U^r/(r-1)!\int_0^1(1-z)^{r-1}g_m^{(r)}(1+zU)\,dz$ and
$r(1-z)^{r-1}$ is a probability density. Finiteness of the required
Gaussian ratio moment justifies expectation over $y$. The price is the
sum of the already evaluable coefficients $c_0,\ldots,c_{r-1}$ and
$\mathbb E R_r$.

For an independent $S$ uniform on $[0,1]$, set
\[
 x(S)=\log(1+HS/K),\quad t=1+ZU,\quad
 \xi=(m+\log t-x(S))/\lambda.
\]
The strike-measure derivative identity gives
\begin{equation}
 \mathbb E R_r={H\over r!}\mathbb E_{Y,Z,S}
       [U^r t^{-r}A_{r,\lambda}(\xi)\phi(\xi)].
 \label{eq:quantum-density-remainder}
\end{equation}
This identity uses a Gaussian density and a finite polynomial, without
a normal-CDF subtraction or division by a small $U$. It does introduce
auxiliary randomness, whose variance and gate costs must be compared
with direct evaluation of $R_r$.

On a certified financial support with $0\leq U\leq U_*$ and $U_*>0$,
choose an exact rational certified upper norm $\overline D_r\geq D_r$ and
define
\begin{equation}
 X_r=(U/U_*)^r t^{-r}A_{r,\lambda}(\xi)\phi(\xi)/\overline D_r,
 \qquad B_r=H\overline D_r U_*^r/r!.
 \label{eq:quantum-remainder-signal}
\end{equation}
Then $|X_r|\leq1$ and the support-restricted correction is $B_r\mathbb E X_r$.
For $U_*=0$ the correction is identically zero and no oracle is needed.
Both $U$ and $m$ are functions of the original classical Gaussian factor;
the method does not introduce quantum asset-price dynamics.

\begin{proposition}[Signed finite comparator encoding]
\label{prop:integral-signed-comparator}
Suppose a reversible value calculation supplies a signed approximation
$\widetilde X$ with pointwise error at most $\eta$ on the retained inputs.
Clamp it to $[-1,1]$ and set
$a=\lfloor 2^q(\widetilde X+1)\rfloor\in\{0,\ldots,2^{q+1}\}$.
The threshold output uses $q+2$ bits to represent the saturated endpoint.
A uniform $(q+1)$-bit register $J$, zero-extended to $q+2$ bits, and the strict comparison flag
$[J<a]$ give success probability $p=\mathbb E[a/2^{q+1}]$.
The decoded signed mean $2p-1$ differs from $\mathbb E X_r$ by at most
$\eta+2^{-q}$.
\end{proposition}
\begin{proof}
Projection onto $[-1,1]$ cannot increase distance from the true $X_r$.
For each integer threshold there are exactly $a$ successful comparison
values, including the zero and saturated endpoint cases. Rounding gives
$0\leq\widetilde X-(a/2^q-1)<2^{-q}$. Taking expectations and using the
pointwise approximation bound proves the result. Original latent and
auxiliary labels are retained, so their orthogonality makes the Born
probability the stated mixture, without interference between unequal labels.
The clean comparison writes into a retained success flag. I uncompute
the value calculation and private comparator work with their coherent
adjoints while excluding that flag-writing action from the reversed body.
The success flag is retained and private work returns to zero.
\end{proof}

For $r=1$, $A_{1,\lambda}=1/\lambda>0$, so $X_1\in[0,1]$.
I may use the direct threshold $a=\lfloor2^q\widetilde X_1\rfloor$,
with $q+1$ threshold bits and a uniform $q$-bit comparison register.
Exactly $a$ outcomes succeed. After projection of the numerical value
onto $[0,1]$, its mean bias is at most $\eta+2^{-q}$ and the monetary
estimation radius is $B_1e_p$. This positive case avoids a signed offset;
it is closely related to the already implemented geometric-control
residual and is an essential opponent rather than a new primitive.
For higher orders I require a proved sign condition before using this
simpler encoding.

The quantum mean estimator acts on this flag. A scalar probability radius
$e_p$ contributes monetary correction radius $2B_r e_p$. Uniform value
error and rounding contribute $B_r(\eta+2^{-q})$. The eventual implementation must charge preparation, controlled
reflections, confidence amplification and reset under the accounting
conventions of Section~\ref{sec:canonical-qpe-comparator}. The statement is
a specified encoding interface; it does not assign zero cost to its
financial value calculation.

\begin{proposition}[Auxiliary randomisation and the competent comparator]
\label{prop:integral-randomisation-variance}
On the retained support, let $\mathcal R=R_r(Y)$ and use the exact
continuous auxiliary laws. Then
\[
 B_r^2\operatorname{Var}(X_r)
 =\operatorname{Var}(\mathcal R)
   +B_r^2\mathbb E[\operatorname{Var}(X_r\mid Y)],\qquad
 \mathbb E X_r^2\leq\mathbb E[(U/U_*)^{2r}].
\]
Thus auxiliary randomisation does not reduce variance against exact
conditional correction evaluation. Any advantage must arise from its
complete evaluator cost, its usable error certificate or a different
proved estimator consequence.
\end{proposition}
\begin{proof}
Equation~\eqref{eq:quantum-density-remainder} holds conditionally on $Y$,
so $B_r\mathbb E[X_r\mid Y]=R_r(Y)$. Expanding the square of
$X_r-\mathbb E X_r$ into its conditional-mean and centred parts gives
the total-variance identity: the cross-term vanishes by conditional
expectation. The squared norm bound follows from $t^{-r}\leq1$ and
$|A_r\phi|\leq\overline D_r$. All variables are bounded on the retained
support, so these expectations exist. A canonical success flag has variance
$p(1-p)$ for $p=(1+\mathbb E X_r)/2$; this is a different observable and
must not be assigned $\operatorname{Var}(X_r)$ without an estimator proof.
\end{proof}

\subsection{An emitted minimum circuit for the auxiliary law}

If $V_1,\ldots,V_r$ are independent uniforms, then
$\Pr(\min_i V_i>z)=(1-z)^r$, so their minimum has the required Beta law.
Prepare $r$ independent $m$-bit registers by $rm$ Hadamards and interpret
each integer $j_i$ as $(j_i+1/2)2^{-m}$. The reversible minimum computes
the minimum integer into a fresh register. For two inputs $a,b$, compute
$c=[a<b]$ and each difference bit $a_i\mathbin\oplus b_i$, then write
$b_i\mathbin\oplus[c(a_i\mathbin\oplus b_i)]$. This is the correct
minimum for both values of $c$, including ties. Iterate the construction,
copy the final result and reverse every private calculation. The original
uniform registers are retained. Linearity of this exact basis permutation
proves its coherent action, including an arbitrary reference system.
Its clean inverse remains available inside amplification and estimation.

The realised constructor passes 10,076 basis tests, including dirty output
registers, and exactly reproduces the histogram
$((N-j)^r-(N-j-1)^r)/N^r$ for tested small registers. An independent audit verifies a further 38,828 dirty-output basis cases
and three full complex Clifford and T replays with a reference system.
It also checks the Hadamard-prepared Born histogram. Forty emitted profiles
cover orders one through eight and widths one, two, four, eight and sixteen.
For $r=2$, $m=8$, the realised minimum circuit has 112 exact Toffolis,
hence 784 T gates under the specified seven-T expansion, plus 16 input
Hadamards. Its 50-wire allocation has 24 retained input/output wires;
private work is zero. These are auxiliary costs only. The inherited
three-layer Toffoli assignment is not used as a new depth claim: I retain
the independently checked four-layer serial and conservative five-layer
caps separately.

Couple each continuous uniform to the midpoint of its containing cell.
Each changes by at most $2^{-m-1}$. The minimum is one-Lipschitz in the
maximum coordinate distance: if all inputs change by at most $h$, both
minima lie within $h$ by comparing each original input with its changed
counterpart. Therefore $|Z-Z_m|\leq2^{-m-1}$. Differentiating the integrand
in equation~\eqref{eq:quantum-beta-remainder} with respect to $z$ gives
$U g_m^{(r+1)}(1+zU)$. The joint norm theorem consequently bounds the
monetary Beta-grid bias by
\begin{equation}
 \beta_Z\leq {H\overline D_{r+1}\over r!}
           \mathbb E[U^{r+1}]\,2^{-m-1}.
 \label{eq:quantum-beta-grid-bias}
\end{equation}
This is a finite-register error bound rather than ideal continuous sampling.

For a midpoint strike register of width $b$, define
$E_r=\sup_z|(A'_{r,\lambda}-zA_{r,\lambda})\phi(z)|$.
It has the independently valid upper bound
$\overline E_r=\sum_{\ell=1}^r|s(r,\ell)|\lambda^{-\ell}\overline{\mathcal M}_\ell$,
using the retained individual Hermite norms. The polynomial recurrence
$A_{r+1,\lambda}=-rA_{r,\lambda}+(A'_{r,\lambda}-zA_{r,\lambda})/\lambda$
also gives $E_r\leq\lambda(D_{r+1}+rD_r)$ by the triangle inequality.
I may take the smaller of these two certified upper bounds; the retained
interface audit verifies all 24 relevant recurrences exactly. Since
$|x'(S)|\leq H/K$, differentiation of the density integrand gives
\begin{equation}
 \beta_S\leq {H^2\over K\lambda r!}\overline E_r
                   \mathbb E[U^r]2^{-b-1}.
 \label{eq:quantum-strike-grid-bias}
\end{equation}
The factor $t^{-r}\leq1$ and the nonnegative moment justify taking
expectations. Gaussian support approximation, coefficient arithmetic,
financial-value precision and state-preparation error are additional
charges, not included in these two auxiliary bounds.

\subsection{Transfer from the Gaussian model to a finite financial register}

I make the retained law explicit. Let $C_L=[-L,L]^p$, $L>0$, and
$a_L=[\Phi(L)-\Phi(-L)]^p$. The conditional law $Y_L=Y\mid C_L$ is a
product of truncated standard normals. Its contribution to the correction
is $a_L B_r\mathbb E[X_r\mid C_L]$, not $B_r\mathbb E[X_r\mid C_L]$.
The omitted correction is bounded by the residual envelope rather than
only the cap's tail probability.

With $b_r=H\overline D_r/r!$ and the finite Gaussian exposure notation of
Section~\ref{sec:taylor-gaussian-tilt}, positivity gives
\begin{equation}
 \beta_{\rm tail}\leq b_r\sum_{|\alpha|=r}{r!\over\alpha!}w^\alpha
 e^{d_\alpha+\|v_\alpha\|^2/2}
 \sum_{j=1}^p\{\Phi(v_{\alpha,j}-L)+\Phi(-v_{\alpha,j}-L)\}.
 \label{eq:quantum-remainder-tail}
\end{equation}
Indeed $U^r\leq T^r$, every multinomial term is nonnegative, and Gaussian
tilting turns its tail expectation into its exponential prefactor times
the probability that a normal vector of mean $v_\alpha$ leaves the box.
The union bound gives the sum of its coordinate tails. This derivation
does not rely on cancellation in a tail moment or assume that the
preconditioned residual is globally bounded by $H$.

Within the box choose certified bounds $U\leq U_*$ and
\[
 B_j^U=\sum_iw_i|s_{i,j}|
          \exp(d_i+L\|s_i\|_1)\geq|\partial_j U|.
\]
Both the exact and auxiliary-discretised correction have coordinate
Lipschitz constants bounded by
\begin{equation}
 L_j^R={H\over r!}\left[
 rU_*^{r-1}B_j^U\overline D_r+
 U_*^r\left({|R_{G,j}|\over\lambda}\overline E_r+
                              B_j^U\overline D_{r+1}\right)\right].
 \label{eq:quantum-remainder-grid-slope}
\end{equation}
Differentiate $U^r g_m^{(r)}(1+zU)/r!$ under its bounded auxiliary
expectation. The derivative of $U^r$ gives the first term; changing $m$
gives the $\overline E_r/\lambda$ term; changing $t$ gives the last term,
using $z\leq1$. For a fixed strike point the same bounds hold, because
its $t$ derivative satisfies
$A_{r+1}= -rA_r+(A_r'-zA_r)/\lambda$.
Thus convex averaging under finite strike and Beta laws preserves these
constants. At $U=0$ the polynomial factor has its ordinary continuous
derivative; no quotient is introduced. The identically zero-support case
is handled before applying this formula.

Prepare each conditional-normal coordinate from a uniform $N$-point
midpoint quantile register. A monotone quantile has total variation $2L$;
on each uniform cell its integrated distance from its midpoint value is
at most half the cell width times the endpoint variation. Summing gives
expected coordinate error at most $L/N$. If the numerical node error is
at most $\eta_Y$, with nodes clamped into the box, the extra distance is
at most $\eta_Y$. The mean financial-grid bias is therefore
\begin{equation}
 \beta_Y\leq a_L\sum_{j=1}^p L_j^R(L/N+\eta_Y).
 \label{eq:quantum-remainder-financial-grid}
\end{equation}
The coupling is coordinatewise under the product truncated law, and a
telescoping replacement of coordinates uses their uniform Lipschitz bounds.

To specify the moment law in the auxiliary certificates, I first replace
the continuous Beta and strike laws while $Y$ still has its continuous
conditional law. Multiplication by $a_L$ replaces each conditional moment
by $\mathbb E[U^r\mathbf1_{C_L}]\leq\mathbb E[U^r]$.
I may therefore use the continuous Gaussian moments in
equations~\eqref{eq:quantum-beta-grid-bias} and
\eqref{eq:quantum-strike-grid-bias} as upper bounds on these unnormalised
auxiliary errors. I then replace $Y_L$ by its finite register using
equation~\eqref{eq:quantum-remainder-financial-grid} for the
already discretised auxiliary function. This order avoids substituting
a continuous moment for a different finite-grid moment without proof.

If coefficient arithmetic has radius $\varepsilon_C$, numerical $a_L$
has error $\eta_a$, value error is $\eta$, and the scalar estimator has
radius $e_p$ at its stated confidence, an undiscounted error envelope is
\begin{equation}
 \varepsilon_C+\beta_{\rm tail}+\beta_Z+\beta_S+\beta_Y
       +B_r\eta_a+a_LB_r(\eta+2^{-q}+2e_p).
 \label{eq:quantum-remainder-full-interface}
\end{equation}
I project the final probability estimate onto $[0,1]$ before decoding.
Projection cannot increase its error from the true probability, and the
decoded signed estimate then lies in $[-1,1]$. This supplies the factor
$B_r$ multiplying $\eta_a$, including the normaliser cross-term. Alternatively, outward interval multiplication propagates
this normaliser directly. Logical Hadamards and the reversible finite maps are assumed exact here;
node approximation is charged by $\eta_Y$ and value approximation by $\eta$.
If an approximate input preparation is used, it requires its own certified
charge. For a success projector $P$, the signed observable $O=2P-I$ has
operator norm one. For normalised states $\psi,\widetilde\psi$ with
$\|\psi-\widetilde\psi\|\leq\delta$, expanding
$\langle\psi,O\psi\rangle-\langle\widetilde\psi,O\widetilde\psi\rangle$
into two terms and applying Cauchy--Schwarz bounds its magnitude by
$2\delta$. Additional correction radius $2a_LB_r\delta$ is therefore
sufficient. No such approximation is
silently assigned zero cost. Discounting and its arithmetic are then
propagated separately. These are sufficient bounds; numerical choices
must demonstrate that their sum fits the requested tolerance. The old
finite-output bias allowance is not transferred automatically to this
different residual.

Before invoking the quantum correction I also test
$|\mathbb E R_r|\leq b_r\mathbb E U^r$. If this whole-correction interval
already meets the remaining tolerance, its midpoint suffices with no
quantum calls. This supplies the same information to the classical and
quantum competitors and prevents claiming a quantum benefit for an
already resolved price.

\subsection{Why a signed Gaussian importance loader is not the default}

I also tested a coherent Gaussian-tilt proposal for an even-power remainder
envelope. Expanding $(T-1)^s\sqrt{p}$ yields a signed finite combination
of translated Gaussian states. Its squared target norm is
$M_{2s}=\mathbb E[(T-1)^{2s}]$, while its coefficient one-norm is $L_s$.
Established LCU gives success $M_{2s}/L_s^2$; this success and the
component precision cannot be omitted \citep{rosenkranz2025,ozols2013}.
The retained toy signal requires 4,416 T gates before amplification,
and 173,288 after its specified amplification. At its tested dispersion,
the latter has ideal coherent-replay success with probability approximately $0.99836$.
An equally informed direct finite-support signal costs 736 T gates and
defeats the candidate by approximately $18.44$ in the signal-only,
correction-weight-adjusted comparison. The complete financial observable
and pricing estimator were not implemented in that toy, so this ratio is
not a complete pricing comparison.

The structural objection has specific assumptions. For this comparison both methods use the same fixed normalised base law
on a fixed certified support, including the same financial correction. Take
centred component exponents $d_i(\delta)+s_i(\delta)\cdot y
=\delta h_i(y)+\delta^2 k_i(y)+O(\delta^3)$ uniformly, with
$\sum_iw_ih_i=\sum_iw_ik_i=0$ and $h_i$ affine. Exponential expansion gives
$U=\delta^2 Q(y)+O(\delta^3)$, where
$Q(y)=\tfrac12\sum_iw_ih_i(y)^2$.
Assume that this base law satisfies $\mathbb E Q^{2s}>0$ and does not depend
on $\delta$. Then $M_{2s}=\Theta(\delta^{4s})$, while the displayed unmerged, labelled expansion's
coefficient one-norm tends to $2^s$. Its state-loading factor is therefore
$\Theta(\delta^{-2s})$, with the common smoothing $\lambda$ fixed.
In the active estimation regime $\varepsilon=o(\delta^{4s})$, its ideal
amplified estimation-call scale is $O(\delta^{2s}/\varepsilon)$.
Direct remainder range on that same fixed support is $O(\delta^{4s})$,
giving the smaller direct call scale $O(\delta^{4s}/\varepsilon)$ before
confidence and evaluator costs. Fixed unequal initial prices or unscaled
drift differences need not satisfy these assumptions. At a fixed tolerance,
sufficiently small corrections enter the zero-query certificate regime;
the asymptotic comparison does not assert fractional physical calls. A positive Gaussian-mixture rejection proposal supplies a further
published-mechanism competitor. Exact merging of identical component states can change this coefficient
norm and requires a separate derivation. These findings reject an unsupported
advantage claim for the tested signed loader and motivate the direct integral
encoding above.

The finite bias certificates currently retained provide $D_{r+1}$ and
the required individual derivative norms for $r\leq8$. A ninth-order
remainder would require new tenth-order and ninth individual norm
certificates before those auxiliary bounds could be used numerically.

The current development is therefore a finance-specific oracle proposal
with a proved correction identity, normalisation and auxiliary circuit.
Taylor expansion, threshold smoothing, comparator encoding and amplitude
estimation are established components. A positive literature contribution
requires the full value oracle, its monetary error certificate and a
competitive comparison with direct residual evaluation, QSP where
compatible, variance-sensitive quantum methods and classical pricers.
The changed research direction permits improving an established quantum
method; it does not turn a new name or an unpriced oracle into a contribution.
